% GENERATED FILE. Edit canonical lessons, then run npm run generate:books.
% canonical-source-hash: fnv1a64:0b09db9d0f63c227
% canonical-lessons: GE-C515-chor, GE-C515-orchester, GE-C515-dirigent, GE-C515-dirigieren, GE-C515-komponieren

\chapter{Chor und Orchester}
\label{ch:ge-chor-und-orchester}
\hlchaptermodality{515}

\begin{chapteropening}
\textbf{By the end of this chapter:} \emph{I can talk about choirs, orchestras and composers.}

Complete the last lesson of chapter 515: I can talk about choirs, orchestras and composers.
\end{chapteropening}

\section[der Chor]{der Chor — choir}

\label{lesson:GE-C515-chor}



\begin{quote}
Before the new one: say the German for verse, then the German for to hum.
\end{quote}

\subsection*{You'll want to know: der Chor}
\textbf{der Chor} — \textquotedblleft{}choir\textquotedblright{}.

\begin{quote}
Meine Mutter singt in einem Chor.
\end{quote}

\begin{quote}
\emph{My mother sings in a choir.}
\end{quote}

Plural: \emph{die Chöre}. \emph{Ch} sounds like \emph{k}.

\begin{cousinweb}
\emph{Chor} is a learned borrowing from Latin \emph{chorus}, from Greek \emph{chorós} \textquotedblleft{}a dance, a band of dancers and singers\textquotedblright{} (Lewis \& Short).
\end{cousinweb}

\begin{grammarlens}[title={What you've built}]
One more word for telling, explaining and describing at B1.
\end{grammarlens}

\subsection*{Your turn}
\begin{itemize}
\raggedright
  \item \emph{Say it:} \emph{der Chor}
  \item \emph{Say it:} \emph{der Chor}, once more
  \item \emph{Say it:} a sentence of your own with \emph{der Chor}
  \item \emph{Recall:} say the German for verse, then the German for to hum, then say \emph{der Chor} again
\end{itemize}

\begin{tcolorbox}[breakable,colback=teal!4,colframe=teal!35!black,title={Before you move on}]
What does der Chor mean? (\textquotedblleft{}Choir\textquotedblright{}.) Say it once more.
\end{tcolorbox}
\section[das Orchester]{das Orchester — orchestra}

\label{lesson:GE-C515-orchester}



\begin{quote}
Before the new one: say the German for to hum, then the German for choir.
\end{quote}

\subsection*{You'll want to know: das Orchester}
\textbf{das Orchester} — \textquotedblleft{}orchestra\textquotedblright{}.

\begin{quote}
Das Orchester hat sechzig Musiker.
\end{quote}

\begin{quote}
\emph{The orchestra has sixty musicians.}
\end{quote}

Plural: \emph{die Orchester}. \emph{Ch} sounds like \emph{k}.

\begin{cousinweb}
\emph{Orchester} comes, through Latin \emph{orchēstra}, from Greek \emph{orch\'{\={e}}stra}; in the Roman theatre it was the space in front of the stage where the senators sat (Lewis \& Short), and the musicians took it over later.
\end{cousinweb}

\begin{grammarlens}[title={What you've built}]
One more word for telling, explaining and describing at B1.
\end{grammarlens}

\subsection*{Your turn}
\begin{itemize}
\raggedright
  \item \emph{Say it:} \emph{das Orchester}
  \item \emph{Say it:} \emph{das Orchester}, once more
  \item \emph{Say it:} a sentence of your own with \emph{das Orchester}
  \item \emph{Recall:} say the German for to hum, then the German for choir, then say \emph{das Orchester} again
\end{itemize}

\begin{tcolorbox}[breakable,colback=teal!4,colframe=teal!35!black,title={Before you move on}]
What does das Orchester mean? (\textquotedblleft{}Orchestra\textquotedblright{}.) Say it once more.
\end{tcolorbox}
\section[der Dirigent]{der Dirigent — conductor (of music)}

\label{lesson:GE-C515-dirigent}



\begin{quote}
Before the new one: say the German for choir, then the German for orchestra.
\end{quote}

\subsection*{You'll want to know: der Dirigent}
\textbf{der Dirigent} — \textquotedblleft{}conductor (of music)\textquotedblright{}.

\begin{quote}
Am Ende verbeugte sich der Dirigent.
\end{quote}

\begin{quote}
\emph{At the end the conductor took a bow.}
\end{quote}

Weak noun: \emph{den Dirigenten}. Plural: \emph{die Dirigenten}.

\begin{cousinweb}
\emph{Dirigent} comes from Latin \emph{dīrigēns} \textquotedblleft{}directing, steering\textquotedblright{} (EtymDB), the present participle of \emph{dīrigere} \textquotedblleft{}to direct\textquotedblright{}.
\end{cousinweb}

\begin{grammarlens}[title={What you've built}]
One more word for telling, explaining and describing at B1.
\end{grammarlens}

\subsection*{Your turn}
\begin{itemize}
\raggedright
  \item \emph{Say it:} \emph{der Dirigent}
  \item \emph{Say it:} \emph{der Dirigent}, once more
  \item \emph{Say it:} a sentence of your own with \emph{der Dirigent}
  \item \emph{Recall:} say the German for choir, then the German for orchestra, then say \emph{der Dirigent} again
\end{itemize}

\begin{tcolorbox}[breakable,colback=teal!4,colframe=teal!35!black,title={Before you move on}]
What does der Dirigent mean? (\textquotedblleft{}Conductor (of music)\textquotedblright{}.) Say it once more.
\end{tcolorbox}
\section[dirigieren]{dirigieren — to conduct (music)}

\label{lesson:GE-C515-dirigieren}



\begin{quote}
Before the new one: say the German for orchestra, then the German for conductor.
\end{quote}

\subsection*{You'll want to know: dirigieren}
\textbf{dirigieren} — \textquotedblleft{}to conduct (music)\textquotedblright{}.

\begin{quote}
Heute dirigiert eine junge Frau das Konzert.
\end{quote}

\begin{quote}
\emph{Today a young woman is conducting the concert.}
\end{quote}

Regular, Perfekt without \emph{ge-}: \emph{hat dirigiert}.

\begin{cousinweb}
\emph{dirigieren} is a learned borrowing from Latin \emph{dīrigere} \textquotedblleft{}to run, to direct\textquotedblright{} (EtymDB).
\end{cousinweb}

\begin{grammarlens}[title={What you've built}]
One more word for telling, explaining and describing at B1.
\end{grammarlens}

\subsection*{Your turn}
\begin{itemize}
\raggedright
  \item \emph{Say it:} \emph{dirigieren}
  \item \emph{Say it:} \emph{dirigieren}, once more
  \item \emph{Say it:} a sentence of your own with \emph{dirigieren}
  \item \emph{Recall:} say the German for orchestra, then the German for conductor, then say \emph{dirigieren} again
\end{itemize}

\begin{tcolorbox}[breakable,colback=teal!4,colframe=teal!35!black,title={Before you move on}]
What does dirigieren mean? (\textquotedblleft{}To conduct (music)\textquotedblright{}.) Say it once more.
\end{tcolorbox}
\section[komponieren]{komponieren — to compose}

\label{lesson:GE-C515-komponieren}



\begin{quote}
Before the new one: say the German for conductor, then the German for to conduct.
\end{quote}

\subsection*{You'll want to know: komponieren}
\textbf{komponieren} — \textquotedblleft{}to compose\textquotedblright{}.

\begin{quote}
Er hat schon als Kind Lieder komponiert.
\end{quote}

\begin{quote}
\emph{He was composing songs even as a child.}
\end{quote}

Regular, Perfekt without \emph{ge-}. The composer is \emph{der Komponist}.

\begin{cousinweb}
\emph{komponieren} is a learned borrowing from Latin \emph{compōnere} \textquotedblleft{}to put together\textquotedblright{}, built on \emph{pōnere} \textquotedblleft{}to put\textquotedblright{} (Lewis \& Short).
\end{cousinweb}

\begin{grammarlens}[title={What you've built}]
One more word for telling, explaining and describing at B1.
\end{grammarlens}

\subsection*{Your turn}
\begin{itemize}
\raggedright
  \item \emph{Say it:} \emph{komponieren}
  \item \emph{Say it:} \emph{komponieren}, once more
  \item \emph{Say it:} a sentence of your own with \emph{komponieren}
  \item \emph{Recall:} say the German for conductor, then the German for to conduct, then say \emph{komponieren} again
\end{itemize}

\begin{tcolorbox}[breakable,colback=teal!4,colframe=teal!35!black,title={Before you move on}]
What does komponieren mean? (\textquotedblleft{}To compose\textquotedblright{}.) Say it once more.
\end{tcolorbox}
